\scanpage{17}
2. Soient $C,\deleted{\unsure{\widetilde C}}C'$ deux catégories \deleted{additives, abéliennes} \inserted{additives}, $S:\deleted{\unsure{A\longrightarrow\widetilde A}}C\longrightarrow C'$ un foncteur additif de $C$ dans $\deleted{\unsure{\widetilde C}}C'$. On suppose que pour \deleted{tout} \deleted{\unsure{$A\in C$}}, $X\in\bar C$, le foncteur $A\mathrel{\unsure{-|\to|\mapsto}}\mathrm{Hom}_{C'}(S(A),X)$ est un foncteur du type \deleted{$\mathrm{Hom}(A,$} $h_B$ (où $B$ dépend de $X$, mais de façon essentiellement fonctorielle, donc on aura
% slanted
\margin{Mettre en évidence le rôle symétrique joué par $C$ et $C'^{o}$}
\[
 \mathrm{Hom}_{\bar C}(S(\deleted{\unsure{\widetilde A}}A),X)\xleftarrow{\sim}\mathrm{Hom}_C(A,T(X)) \tag*{(1)}
\]
où $T:\deleted{\unsure{\widetilde C}}C'\longrightarrow C$ est un foncteur covariant exact à gauche de $\deleted{\unsure{\widetilde C|C^{o}}}C'$ dans $C$]. Cette condition est satisfaite notamment si $C$ et $\deleted{\unsure{\widetilde C}}C'$ satisfont à $\mathrm{AB}\unsure{(1,|1,|(}3,4,5)$, si $C$ a un générateur, et si de plus \deleted{\unsure{$A\longmapsto\widetilde A$}} $S$ est un foncteur exact à droite (donc $A\mathrel{-}\mathrm{Hom}_C(\deleted{\unsure{\widetilde A}}S(A),X)$ est exact à gauche). \deleted{En particulier,}

En particulier, si $X=S(\deleted{\unsure{\widetilde B}}B)$, on aura
\[
 \mathrm{Hom}_C(A,B)\longrightarrow\mathrm{Hom}_{C'}(\deleted{\unsure{\widetilde A}}S(A),\deleted{\unsure{\widetilde B}}S(B))\xleftarrow{\sim}\mathrm{Hom}_C(A,TS(B)) \tag*{(2)}
\]
Donc faisant $A=B$, on trouve un élément canonique
\[
 \varphi(A)\in\mathrm{Hom}(A,TS(\deleted{\unsure{\widetilde A}}A))\qquad\text{i.e.} \tag*{(3)}
\]
\[
 \varphi:\mathrm{id}_C\longrightarrow TS \tag*{(3 bis)}
\]
D'autre part, on tire de (1), faisant $A=T(Y)$
\[
 \deleted{\mathrm{Hom}_{\widetilde C}(\widetilde A,X)\xrightarrow{\sim}\mathrm{Hom}_C(A,T(X))\longrightarrow\mathrm{Hom}_{\widetilde C}(\widetilde A,\widetilde{T(X)})} \tag*{\deleted{(4)}}
\]
\[
 \mathrm{Hom}_{C'}(\unsure{TS|ST}(\deleted{\unsure{\widetilde Y}}Y),X)\simeq\mathrm{Hom}_C(T(Y),T(X)) \tag*{(4)}
\]
et en particulier, faisant \deleted{\unsure{A}} $X=Y$, on trouve un élément canonique
\[
 \psi(X)\in\mathrm{Hom}_{C'}(ST\deleted{\unsure{\widetilde{\phantom X}}}(X),X)\simeq \tag*{(5)}
\]
\[
 \psi:ST\longrightarrow\deleted{\unsure{\varphi}}\mathrm{id}_{\widetilde C} \tag*{(5 bis)}
\]
% margin note written sideways
\margin{C'est la donnée des $\psi(X)$ qui précise l'homomorphisme (1), puisqu'on a $\mathrm{Hom}_C(A,T(X))\longrightarrow\mathrm{Hom}_{\widetilde C}(\widetilde A,\widetilde{T(X)})\longrightarrow\mathrm{Hom}_{\widetilde C}(A,X)$}
% archive sheet number: 16
